\begin{figure}
			\resizebox{\textwidth}{!}{%
				\begin{tikzpicture}[
					scale=0.62,
					asta/.style={line width=1.5pt, line cap=round, line join=round, black},
					astaG/.style={line width=1.5pt, line cap=round, line join=round, gray!60},
					nodo/.style={circle, fill=white, draw=black,
						inner sep=0pt, minimum size=6pt, line width=1.2pt},
					nodoA/.style={circle, fill=white, draw=black,
						inner sep=0pt, minimum size=6pt, line width=1.2pt,
						dashed},
					forza/.style={-{Stealth[length=5pt,width=4pt]},
						line width=1.2pt, black},
					lab/.style={font=\normalsize},
					]
					
					\def\SEP{13.0}
					
					%%── PANNELLO SINISTRO: Caso 1, H_B=1 ────────────────────────────
					\begin{scope}[xshift=0cm]
						\coordinate (A) at (0,0);
						\coordinate (B) at (8,0);
						\coordinate (C) at (4,3);
						\coordinate (D) at (4,0);
						
						\draw[asta](A)--(C);
						\draw[asta](C)--(B);
						\draw[asta](A)--(D);
						\draw[asta](D)--(B);
						\draw[asta](C)--(D);
						
						%% cerniera A
						\draw[line width=1.2pt](A)--++(-0.4,-0.65)--++(0.8,0)--cycle;
						\draw[line width=1.2pt](-0.55,-0.65)--(0.55,-0.65);
						\foreach \x in {-0.45,-0.25,-0.05,0.15,0.35}
						\draw[thin](\x,-0.65)--++(-.18,-.22);
						
						%% carrello B
						\draw[line width=1.2pt](B)--++(-0.4,-0.50)--++(0.8,0)--cycle;
						\draw[line width=1.2pt]($(B)+(-0.50,-0.50)$)--($(B)+(0.50,-0.50)$);
						\draw[line width=1.2pt]($(B)+(-0.25,-0.65)$)circle(3.5pt);
						\draw[line width=1.2pt]($(B)+( 0.25,-0.65)$)circle(3.5pt);
						\draw[line width=1.2pt]($(B)+(-0.55,-0.82)$)--($(B)+(0.55,-0.82)$);
						\foreach \x in {-0.45,-0.25,-0.05,0.15,0.35}
						\draw[thin]($(B)+(\x,-0.82)$)--++(-.18,-.22);
						
						\node[nodo]at(A){}; \node[nodo]at(B){};
						\node[nodo]at(C){}; \node[nodo]at(D){};
						
						%% H_B=1 verso destra in B
						\draw[forza](B)--++(1.20,0);
						\node[right,lab]at($(B)+(1.20,0)$){$H_B=1$};
						
						%% H_A=-1 verso sinistra in A (reazione)
						\draw[forza](A)--++(-1.20,0);
						\node[below,lab]at($(A)+(-1.20,0)$){$1$};
						
						\node[above left,lab]at(A){$A$};
						\node[above right,lab]at(B){$B$};
						\node[above right,lab]at(C){$C$};
						\node[above left,lab]at(D){$D$};
						
						\node[font=\normalsize]at(4,-1.60){$(a)$};
					\end{scope}
					
					%%── PANNELLO DESTRO: Caso 2, N_AB=1 ────────────────────────────
					\begin{scope}[xshift=\SEP cm]
						\coordinate (A) at (0,0);
						\coordinate (B) at (8,0);
						\coordinate (C) at (4,3);
						\coordinate (D) at (4,0);
						%% asta AB traslata verso il basso
						\coordinate (Aab) at (0,-0.30);
						\coordinate (Bab) at (8,-0.30);
						
						%% struttura senza AB (completamente scarica)
						\draw[asta](A)--(C);
						\draw[asta](C)--(B);
						\draw[asta](A)--(D);
						\draw[asta](D)--(B);
						\draw[asta](C)--(D);
						
						%% asta AB traslata verso il basso (grigia)
						\draw[astaG](Aab)--(Bab);
						
						%% cerniera A
						\draw[line width=1.2pt](A)--++(-0.4,-0.65)--++(0.8,0)--cycle;
						\draw[line width=1.2pt](-0.55,-0.65)--(0.55,-0.65);
						\foreach \x in {-0.45,-0.25,-0.05,0.15,0.35}
						\draw[thin](\x,-0.65)--++(-.18,-.22);
						
						%% carrello B
						\draw[line width=1.2pt](B)--++(-0.4,-0.50)--++(0.8,0)--cycle;
						\draw[line width=1.2pt]($(B)+(-0.50,-0.50)$)--($(B)+(0.50,-0.50)$);
						\draw[line width=1.2pt]($(B)+(-0.25,-0.65)$)circle(3.5pt);
						\draw[line width=1.2pt]($(B)+( 0.25,-0.65)$)circle(3.5pt);
						\draw[line width=1.2pt]($(B)+(-0.55,-0.82)$)--($(B)+(0.55,-0.82)$);
						\foreach \x in {-0.45,-0.25,-0.05,0.15,0.35}
						\draw[thin]($(B)+(\x,-0.82)$)--++(-.18,-.22);
						
						%% cerchi aperti agli estremi dell'asta AB (sconnessa)
						\draw[dashed, line width=0.8pt](Aab)circle(4pt);
						\draw[dashed, line width=0.8pt](Bab)circle(4pt);
						
						\node[nodo]at(A){}; \node[nodo]at(B){};
						\node[nodo]at(C){}; \node[nodo]at(D){};
						
						%% forza unitaria verso destra all'estremo destro (sconnesso)
						\draw[forza](Bab)--++(1.20,0);
						\node[right,lab]at($(Bab)+(1.20,0)$){$N_{AB}=1$};
						
						%% reazione del vincolo in A sull'asta AB: coda in Aab, punta a sinistra
						\draw[forza](Aab)--++(-1.20,0);
						\node[below,lab]at($(Aab)+(-1.20,0)$){$1$};
						
						\node[above left,lab]at(A){$A$};
						\node[above right,lab]at(B){$B$};
						\node[above right,lab]at(C){$C$};
						\node[above left,lab]at(D){$D$};
						
						\node[font=\normalsize]at(4,-1.60){$(b)$};
					\end{scope}
					
				\end{tikzpicture}
			}
			\caption{Autosoluzioni delle due varianti iperstatiche.
				$(a)$~Caso~1 (iperstaticità esterna): forza unitaria
				$H_B = 1$ applicata in~$B$ sul sistema principale;
				la reazione del vincolo in~$A$ equilibra la forza
				unitaria direttamente.
				$(b)$~Caso~2 (iperstaticità interna): l'asta~$AB$
				(in grigio) è sconnessa in~$B$ e soggetta allo
				sforzo unitario $N_{AB} = 1$; la reazione del
				vincolo in~$A$ si applica alla sola asta~$AB$,
				lasciando la struttura rimanente completamente
				scarica.}
			\label{fig:capriata_autosoluzioni}
		\end{figure}
